\section{M4 --- Detecção e Limiares}
\label{sec:m4}

\subsection{Erro de reconstrução}
A detecção opera sobre o erro de reconstrução por janela --- o erro médio absoluto (MAE)
entre a janela de entrada e sua reconstrução. Optou-se pelo MAE por sua maior robustez a
\emph{outliers} isolados em relação ao MSE (este último disponível como opção). Uma janela
é classificada como anômala quando seu erro supera o limiar.

\subsection{Limiar estático}
O limiar estático é o percentil 95 do erro de reconstrução \emph{do treino}. Por ser
calculado sobre o período de normalidade (e não sobre o teste), não incorre em vazamento.
É simples e estável, mas, por definição, fixo: não se adapta a mudanças no regime de
volatilidade do ativo.

\subsection{Limiar dinâmico}
O limiar dinâmico é o percentil 95 calculado em uma janela móvel \textbf{causal} de
$\texttt{dynamic\_window}=252$ observações ($\approx$1 ano útil). A causalidade é garantida
excluindo-se o próprio ponto da janela (deslocamento de uma posição), de modo que uma
anomalia não infle o limiar que a avalia; os primeiros pontos, com janela incompleta, não
são marcados. Esse esquema constitui uma das contribuições metodológicas do trabalho.

\subsection{Comparação empírica}
A Tabela~\ref{tab:thresholds} compara, no conjunto de teste (2020--2024), a fração de
janelas marcadas por cada esquema. O resultado evidencia o comportamento esperado:
\begin{itemize}[noitemsep]
  \item Em ativos de regime relativamente estável (PETR4, VALE3), os dois limiares quase
        coincidem ($\approx$3\%).
  \item Sob mudança abrupta de regime, o limiar estático \textbf{satura}: no AMER3, todo o
        período pós-fraude/recuperação judicial faz o erro permanecer elevado, e o estático
        marca $\approx$27\% das janelas --- perdendo poder discriminativo. O limiar dinâmico,
        ao se recalibrar à volatilidade local, reduz a marcação para $\approx$14\% sem usar
        informação do futuro. Comportamento análogo no ITUB4 (11\% $\to$ 5\%).
\end{itemize}

\begin{table}[h]
\centering
\caption{Fração de janelas de teste marcadas como anômalas, por limiar.}
\label{tab:thresholds}
\begin{tabular}{lccc}
\toprule
\textbf{Ticker} & \textbf{Limiar estático} & \textbf{Estático (\%)} & \textbf{Dinâmico (\%)} \\
\midrule
PETR4.SA & 0{,}124 & 2{,}9 & 2{,}9 \\
VALE3.SA & 0{,}082 & 2{,}7 & 3{,}4 \\
AMER3.SA & 0{,}077 & 26{,}7 & 14{,}3 \\
ITUB4.SA & 0{,}088 & 11{,}0 & 4{,}9 \\
\bottomrule
\end{tabular}
\end{table}

\begin{figure}[h]
\centering
\includegraphics[width=\textwidth]{m4_erro_limiares}
\caption{Erro de reconstrução no teste (2020--2024) com limiar estático (p95, tracejado
vermelho) e dinâmico (verde). O dinâmico acompanha o regime de volatilidade; no AMER3 o
estático satura no período pós-fraude.}
\label{fig:m4-limiares}
\end{figure}

\paragraph{Observações metodológicas.} O percentil 95 implica, por construção, uma taxa-base
de $\approx$5\% de marcações --- a considerar na leitura das métricas (M5). Além disso, como
as janelas se sobrepõem, uma única anomalia gera uma rajada de detecções contíguas; antes de
confrontar com eventos reais (M6) ou calcular métricas (M5), detecções adjacentes devem ser
agrupadas em ``eventos''.

\subsection{Sensibilidade da janela dinâmica}
\label{sec:m4-sensibilidade}
O tamanho \texttt{dynamic\_window} era provisório; a Tabela~\ref{tab:sensibilidade} fecha a
questão variando-o em $\{126, 252, 504\}$. O comportamento é monotônico: a janela curta (126)
é reativa --- segue o ruído local e infla a marcação ($\approx$8\% nos ativos estáveis),
enquanto a janela longa (504) super-suaviza, aproximando-se de um limiar global e colapsando a
marcação a $\approx$0\% --- perdendo a capacidade de sinalizar. O valor intermediário
\textbf{252} produz taxas próximas da base-rate alvo ($\approx$5\% do p95) sem oscilar nem
colapsar, e é confirmado como escolha definitiva (ADR-0005).

\begin{table}[h]
\centering
\caption{Fração de janelas de teste marcadas pelo limiar dinâmico, por tamanho de janela.}
\label{tab:sensibilidade}
\begin{tabular}{lccc}
\toprule
\textbf{Ticker} & \textbf{126} & \textbf{252} & \textbf{504} \\
\midrule
PETR4.SA & 8{,}5\%  & 2{,}9\%  & 0{,}0\% \\
VALE3.SA & 8{,}2\%  & 3{,}4\%  & 0{,}3\% \\
AMER3.SA & 15{,}2\% & 14{,}3\% & 5{,}7\% \\
ITUB4.SA & 6{,}3\%  & 4{,}9\%  & 0{,}0\% \\
\bottomrule
\end{tabular}
\end{table}
